%%% FIELD stmt-multihomogeneous-bezout
Let \(f_1,\ldots,f_q\) be a square polynomial system whose variables are partitioned into blocks of sizes \(r_1,\ldots,r_b\), with \(\sum_i r_i=q\), and let the degree of \(f_\ell\) in block \(i\) be at most \(\alpha_{\ell i}\). The number of isolated affine common zeros, counted with multiplicity, is at most the coefficient of \(\prod_i X_i^{r_i}\) in
\[
\prod_{\ell=1}^{q}\left(\sum_{i=1}^{b}\alpha_{\ell i}X_i\right).
\]
%%% FIELD stmt-generic-rur-specialization
For a generically zero-dimensional polynomial system with parameter variables, a generic rational univariate representation based on a generically separating linear form specializes to a rational univariate representation, with multiplicities preserved, away from the finite union of the algebraic loci where dimension, a leading coefficient, a point at infinity, or separation degenerates.
%%% FIELD stmt-semialgebraic-stratification
Every real semialgebraic set is a finite disjoint union of smooth semialgebraic manifolds satisfying the frontier condition; in particular, every point belongs to a smooth stratum and each stratum is locally Euclidean.
%%% FIELD stmt-schwartz-zippel
If a nonzero polynomial has total degree at most \(d\), evaluation at independent uniform points of a finite set \(S\) vanishes with probability at most \(d/|S|\).
%%% FIELD edit-orientation-statement
For every \((G,T,e)\in\mathcal G_{n,k}\) with \(d_e(G)<\infty\),
\[
d_e(G)\le 2^{|D|-n+2}\le2^{k+1}.
\]
The first inequality holds even when the generic full parameter fiber is positive dimensional.
%%% FIELD edit-exceptional-statement
Assume \(d_e(G)<\infty\). With probability at least \(1-\delta\), \(h_{G,e}\) is not the zero polynomial and \(\mathcal E_{G,e}\) is proper. For every complex covariance specialization \(\Sigma\notin\mathcal E_{G,e}\) in the model image,
\[
\{z\in\mathbb C:p_e(\Sigma,z)=0\}=\mathcal Z_e(G,\Sigma).
\]
This equality uses the specialized prime incidence ideal \(I_G\) and remains valid when its full \(\Lambda\)-fiber is positive dimensional. The circuits in \(\mathcal Q_{G,e}\) certify that every coordinate-preserving slice meets every generic finite branch, intermediate poles and identically zero-over-zero components are retained or lifted back to \(I_G\), the specialized leading coefficient does not drop, and no finite root is exchanged with a root at projective infinity. The roots are distinct whenever the squarefreeness circuit is nonzero.
%%% FIELD edit-regular-cell-statement
Assume \(d_e(G)<\infty\) and the certificate event on which \(h_{G,e}\) is nonzero. The real supported source set
\[
\{(\Lambda,\Omega):\Lambda\in\mathbb R^D,\ \Omega\succ0,\ \Omega_{ij}=0\text{ if }i\leftrightarrow j\notin B\}
\]
is Euclidean open in its support space and Zariski dense in the complex source space. Since \(\mathcal E_{G,e}\) is proper on the model image, there is a real supported source point \((\Lambda_0,\Omega_0)\) whose covariance \(\Sigma_0\) lies in \(\mathcal M_G\setminus\mathcal E_{G,e}\). A smooth semialgebraic stratification refined by the real-lifting sign conditions contains \(\Sigma_0\) in a connected cell \(\mathcal C_{G,e}\) on which real-extension status is constant, and a sufficiently small closed chart ball gives a nonempty compact \(\mathcal K_G\subset\mathcal C_{G,e}\) with nonempty relative interior. For every \(\Sigma\in\mathcal K_G\), \(\mathcal R_e(G,\Sigma)\ne\varnothing\), because membership in \(\mathcal M_G\) supplies at least one complete real supported \(\Lambda\).
%%% FIELD proof-incidence-localization
Put \(R=I_n-\Lambda\). By acyclicity, \(\Lambda^n=0\) after a topological ordering, and hence \(R^{-1}=I_n+\Lambda+\cdots+\Lambda^{n-1}\). Thus the covariance parametrization is polynomial.

Let \(J=(F_{ij}:i\ne j,\ i\leftrightarrow j\notin B)\). In \(R_{\mathrm{inc}}/J\), set
\[
\overline\Omega=R(\overline\Lambda)^{\mathsf T}\overline\Sigma R(\overline\Lambda).
\]
For a missing pair, direct matrix multiplication gives
\[
\overline\Omega_{ij}=\sigma_{ij}
-\sum_{p\in\operatorname{pa}_D(i)}\lambda_{pi}\sigma_{pj}
-\sum_{q\in\operatorname{pa}_D(j)}\lambda_{qj}\sigma_{iq}
+\sum_{p,q}\lambda_{pi}\lambda_{qj}\sigma_{pq}=F_{ij}=0.
\]
Consequently the supported entries of \(\overline\Omega\), together with \(\overline\Lambda\), define a homomorphism \(\theta:R_{\mathrm{src}}\to R_{\mathrm{inc}}/J\). The map induced by \(\psi_G\) is its inverse, because
\[
R(\Lambda)^{\mathsf T}\{R(\Lambda)^{-\mathsf T}\Omega R(\Lambda)^{-1}\}R(\Lambda)=\Omega
\]
and \(R(\overline\Lambda)^{-\mathsf T}\overline\Omega R(\overline\Lambda)^{-1}=\overline\Sigma\). Hence
\[
R_{\mathrm{inc}}/J\cong R_{\mathrm{src}},\qquad J=\ker\psi_G=I_G.
\]
This proves surjectivity, primality, and the quotient assertion. Restriction gives \(P_G=\ker\phi_G\) and \(R_G\cong\operatorname{im}\phi_G\).

The source ring is a domain. Every member of \(S_G\) has nonzero image in it, so \(I_G:S_G^\infty=I_G\) and \(A_K\) is a domain. Moreover \(S_G^{-1}R_G=K_G\), while localization leaves the fraction field unchanged; therefore \(\operatorname{Frac}(A_K)=L\). No directed-coordinate expression was inverted.

Suppose \(\lambda_e\) is algebraic and put \(E=K_G(\lambda_e)\). Bezout's identity modulo its irreducible minimal polynomial shows that every nonzero element of \(K_G[\lambda_e]\) is invertible there, so \(E=K_G[\lambda_e]\subset A_K\). For every \(K_G\)-embedding \(\tau:E\to\overline{K_G}\),
\[
A_K\otimes_{E,\tau}\overline{K_G}\ne0
\]
by faithful flatness. This finitely generated algebra has a geometric point with coordinate \(\tau(\lambda_e)\). Conversely, every geometric point restricts to such an embedding. Thus the coordinate image is exactly the roots of the minimal polynomial, irrespective of the dimension of the full fiber. If \(\lambda_e\) is transcendental, \(K_G[t]\to A_K\), \(t\mapsto\lambda_e\), is injective; the coordinate morphism is dominant, so its image is Zariski dense and infinite.

Finally, the nonzero homogenizing chart dehomogenizes exactly to the affine incidence scheme; the complementary hyperplane consists exactly of projective-infinity points. A denominator zero or leading-coefficient drop is the zero set of a nonzero covariance element. Recording these loci therefore preserves the generic affine coordinate image.
%%% FIELD proof-orientation-upper-bound
Assume \(\lambda_e\) is algebraic. Choose \(U\subseteq\lambda_D\setminus\{\lambda_e\}\) maximal algebraically independent over \(K_G\), put \(F=K_G(U)\), and let \(y\) be the remaining \(q=|D|-|U|\) directed coordinates. Then \(L/F\) is finite separable. Pure transcendence makes \(K_G(U)\) linearly disjoint from \(K_G(\lambda_e)\), whence
\[
[F(\lambda_e):F]=[K_G(\lambda_e):K_G]=d_e(G).
\]
Since \(\Omega_{L/F}=0\) and the residuals generate the prime incidence ideal, their Jacobian in \(y\) has rank \(q\) in \(L\). Select a nonzero \(q\)-square minor. Every \(F\)-embedding of \(F(\lambda_e)\) extends to \(L\), and every conjugate sends this minor to a nonzero one. Thus each of the \(d_e(G)\) queried conjugates occurs at an isolated zero of the selected subsystem.

Partition \(y\) by child vertex. A nonempty block at \(v\) has size \(r_v\le\operatorname{indeg}_D(v)\). A selected residual is unary in one block or bilinear across two blocks. In a connected equation-block component, let \(b,u,m_H,q_C\) denote the numbers of blocks, unary equations, binary equations, and variables. A nonzero determinant term matches rows to incident variable columns, so \(u+m_H=q_C\). For \(b>1\), the binary equations form a connected simple graph \(H\). By \(\mathrm{lem:multihomogeneous\mbox{-}bezout\mbox{-}isolated\mbox{-}bound}\), the isolated-zero bound is the coefficient of \(\prod_vX_v^{r_v}\) in
\[
\left(\prod_{\ell=1}^{u}X_{v(\ell)}\right)
\left(\prod_{\{v,w\}\in E(H)}(X_v+X_w)\right).
\]
This coefficient counts orientations of \(H\) with prescribed indegrees after the unary exponents are removed. Once the \(m_H-b+1\) edges outside a fixed spanning tree are oriented, leaf removal forces at most one orientation of the tree. Hence the coefficient is at most
\[
2^{m_H-b+1}=2^{q_C-u-b+1}
\le2^{\sum_{v\in C}(\operatorname{indeg}_D(v)-1)+1}.
\]
For \(b=1\), the bound is one and the same inequality holds. Applying it only to the component containing the query avoids multiplying unrelated joint-fiber choices. Since every nonroot has an incoming tree edge,
\[
d_e(G)\le2^{\sum_{v\ne\mathrm{root}}(\operatorname{indeg}_D(v)-1)+1}
=2^{|D|-n+2}\le2^{k+1},
\]
where the last step uses \(|D\setminus T|\le k\).
%%% FIELD proof-large-size-frontier
Set \(d=k+1\) and \(n_0=2d+2\). The family has \(3d\) arrows. Its displayed \(T\) has \(2d+1=n_0-1\) arrows and is a rooted spanning arborescence; precisely \(d-1=k\) arrows lie outside it.

Choose positive rational \(h_i\), let \(H_0=\operatorname{diag}(h_i)\), and choose rational \(b_i\) with \(b_i^2>4h_i\). Take a rational positive-definite \(H\) sufficiently close to \(H_0\), with the same diagonal, and prescribe
\[
\sigma_{00}=1,\quad\Sigma_{0P}=\Sigma_{0Q}=0,\quad\sigma_{0h}=1,\quad
\Sigma_{PP}=\Sigma_{QP}=H,\quad\Sigma_{Ph}=0,\quad\Sigma_{Qh}=b.
\]
Choosing \(\Sigma_{QQ}=cI_d\) and then \(\sigma_{hh}=M\) sufficiently large makes \(\Sigma\succ0\) by two Schur complements.

Write \(a_i=\lambda_{0P_i}\), \(c_i=\lambda_{0Q_i}\), \(z_i=\lambda_{P_ih}\), and \(x=-Hz\). The missing \(P_i\)--\(h\) and \(Q_i\)--\(h\) equations force
\[
a_i=x_i,\qquad c_i=x_i+b_i.
\]
The missing \(P_i\)--\(Q_i\) equation is
\[
x_i^2+b_ix_i+h_i=0.
\]
Thus there are exactly \(2^d\) real solutions. Eliminating the two linear groups shows that the residual Jacobian is nonsingular when
\[
\det(H)\prod_i(2x_i+b_i)\ne0,
\]
which holds at all of them.

At diagonal \(H_0\), branches with different first roots have different \(z_1=-(H^{-1}x)_1\). If two branches have the same first root but differ at \(j>1\), perturb \(H_{1j}=H_{j1}\) by \(s\). Since the diagonal and thus \(x\) stay fixed,
\[
\left.\frac{d}{ds}\{z_1(x)-z_1(x')\}\right|_{s=0}
=\frac{x_j-x'_j}{h_1h_j}\ne0.
\]
Each collision is therefore a proper algebraic condition. A small rational perturbation outside their finite union preserves positivity and simplicity and separates all \(2^d\) queried values.

The invertible residual Jacobian recursively solves the formal power-series coefficients of a local inverse over the covariance completion. The \(2^d\) simple points hence give \(2^d\) distinct formal branches with distinct \(z_1\)-constants. Any generic annihilating polynomial for \(z_1\) has all these roots, so its degree is at least \(2^d\); the orientation bound supplies the reverse inequality.

Appending a fully bidirected child of \(0\) introduces no missing residual. Birationally its covariance row and diagonal, together with its free arrow, replace its supported noise coordinates. The new covariance coordinates and arrow are purely transcendental over the old fields, so the old minimal-polynomial degree is unchanged. Therefore, for every \(n\ge2k+4\),
\[
d_{e^\star}(G^\star_{k,n})=2^{k+1},\qquad D_k(n)=2^{k+1}.
\]
%%% FIELD proof-tree-maxima
For a rooted directed tree, write \(x_v\) for the unique coefficient entering nonroot \(v\). A missing root residual is linear in one \(x_v\); a missing nonroot pair is bilinear in \(x_v,x_w\). A rank-two equation gives an invertible Möbius relation. Along each rank-two component all variables are Möbius functions of one seed \(t\), and a nonidentity cycle gives \(t=M(t)\), a polynomial of degree at most two after clearing its denominator. A rank-one residual factors into affine factors. Direct recovery by \(\Omega=(I_n-\Lambda)^{\mathsf T}\Sigma(I_n-\Lambda)\) identifies the residual quotient with the polynomial source ring, so it is a domain; therefore one factor vanishes identically and rationally anchors an endpoint. Thus a finite tree coordinate has degree at most two.

For \(n=2\), absence of the sole residual leaves the query transcendental, while its presence solves the coefficient linearly. For \(n=3\), the rooted shape is a star or chain. There are two coefficients, root residuals are linear anchors, and only one residual can involve both. Without an anchor it leaves a one-dimensional fiber; with an anchor it solves the other coefficient linearly. Unconfounded examples attain degree one. Hence \(D_0(2)=D_0(3)=1\).

For \(n=4\), take the rooted star, allow the three root--child bows, and omit all child--child bidirected pairs. At
\[
(a,b,c)=(1,2,3),\qquad
\Omega=\begin{pmatrix}
1&1/5&-1/7&1/11\\
1/5&10&0&0\\
-1/7&0&10&0\\
1/11&0&0&10
\end{pmatrix},
\]
\(\Omega\succ0\) follows from \(2|u_iu_j|\le u_i^2+u_j^2\) and strict diagonal dominance. Direct multiplication gives
\[
\Sigma=\begin{pmatrix}
1&6/5&13/7&34/11\\
6/5&57/5&79/35&203/55\\
13/7&79/35&94/7&443/77\\
34/11&203/55&443/77&215/11
\end{pmatrix}.
\]
The three residuals have lexicographic reduction
\[
5a+7b-19=0,\qquad5a-11c+28=0,\qquad(a-1)(5a-7)=0.
\]
The lifts are \((1,2,3)\) and \((7/5,12/7,35/11)\). Their Jacobian determinants are \(-2/385\) and \(2/385\), and their Möbius denominators are \(-1/5\) and \(1/5\). Thus two simple, pole-free formal branches exist, proving generic degree two. Fully bidirected root-children preserve the degree by purely transcendental extension. Hence \(D_0(n)=2\) for \(n\ge4\).
%%% FIELD proof-quadratic-core
Let \(S=\{v:\operatorname{indeg}_D(v)\ge2\}\), \(s=|S|\), and expose all \(r=\sum_{v\in S}\operatorname{indeg}_D(v)\) incoming coefficients there. Since \(T\) contributes one arrow to every nonroot,
\[
k_{\mathrm{act}}=|D\setminus T|
=\sum_{v\ne\mathrm{root}}(\operatorname{indeg}_D(v)-1)\le k.
\]
Every member of \(S\) contributes at least one, so \(s\le k_{\mathrm{act}}\) and
\[
r=\sum_{v\in S}(\operatorname{indeg}_D(v)-1)+s\le2k.
\]

Every other nonroot has a scalar coefficient \(y_v\). A rank-two ordinary--ordinary residual gives a Möbius relation. A spanning forest writes \(y_v=M_v(t_C)\); each extra component edge yields a quadratic equation in \(t_C\). A rank-one residual factors into affine factors. Since \(A_K\) is a domain, one factor vanishes identically; source pullback testing selects that factor and supplies a rational anchor. Rank-zero and identically zero-over-zero cases are recorded without division.

After propagation, an ordinary--exposed equation is \(A(z)t_C+B(z)=0\), with \(A,B\) affine in one exposed block. If \(A\ne0\) in the source field, substitute \(t_C=-B/A\) and record \(A\). Two such expressions give the quadratic compatibility equation \(A_1B_2-A_2B_1=0\); a cycle equation remains quadratic after clearing the linear denominator. If \(A=0\), test and retain the affine identity \(B=0\). Exposed--exposed residuals are bilinear. Hence all core equations have degree at most two in at most \(2k\) exposed coefficients and rational-circuit coefficients in \(K_G\).

For every recorded nonzero \(g\in A_K\) and \(E=K_G(\lambda_e)\subset A_K\),
\[
(A_K[1/g])\otimes_{E,\tau}\overline{K_G}\ne0
\]
for every embedding \(\tau\), by faithful flatness. Thus localization preserves every algebraic queried conjugate; for a transcendental query it preserves a dense coordinate image. Complementary pole divisors are checked in the original incidence ideal. An autonomous ordinary component is either free or governed by one quadratic seed, and components not containing the query cannot multiply its coordinate degree. This proves exact preservation of the finite or infinite queried projection.
%%% FIELD proof-exceptional-set
Condition on all identity tests of \(\mathsf A\) being correct. By \(\mathrm{thm:fpt\mbox{-}minimal\mbox{-}polynomial}\), \(p_e\) is the true monic minimal polynomial and every circuit in \(\mathcal Q_{G,e}\) was retained only after its source pullback was found nonzero. Its class in the domain \(R_G\) is nonzero, so their finite product \(h_{G,e}\) is nonzero in \(R_G\). Hence \(\mathcal E_{G,e}\) is proper on the model image.

The relation \(p_e(\lambda_e)=0\) holds in \(A_K\). Clearing covariance denominators yields \(a\in S_G\) and incidence polynomials \(c_{ij}\) with
\[
a(\Sigma)p_e(\Sigma,\lambda_e)
=\sum_{i\leftrightarrow j\notin B}c_{ij}(\Sigma,\Lambda)F_{ij}(\Lambda,\Sigma).
\]
The factors of \(a\) occur in \(\mathcal Q_{G,e}\). Thus outside \(V(h_{G,e})\), every specialized incidence solution gives a root of \(p_e\), proving \(\mathcal Z_e(G,\Sigma)\subseteq V(p_e(\Sigma,\cdot))\).

Conversely put \(E=K_G(\lambda_e)\). For each embedding \(\tau:E\to\overline{K_G}\), the algebra \(A_K\otimes_{E,\tau}\overline{K_G}\) is nonzero. A transcendence basis \(U\) is independent over \(E\), so each conjugate component maps dominantly to \(U\)-space. Their finitely many open images have a common nonempty open intersection. A common coordinate-preserving slice therefore meets every queried branch and is zero dimensional. By \(\mathrm{lem:generic\mbox{-}rur\mbox{-}specialization}\), its rational univariate representation specializes without loss away from its denominator, leading-form, infinity, and separation loci. The returned lifting identities verify by substitution that all RUR points satisfy every original \(F_{ij}\). These finitely many factors are included in \(\mathcal Q_{G,e}\); hence every specialized root lifts to the original incidence scheme outside \(V(h_{G,e})\). This proves the reverse inclusion even for a positive-dimensional unsliced fiber.

The homogenizing leading-form factor separates affine roots from infinity, the pole identities restore zero-over-zero cases, and the specialized coefficient factors prevent degree drops. Characteristic zero makes \(p_e\) separable; nonvanishing of its cleared discriminant, the squarefreeness circuit, makes all specialized roots distinct.

If the tested pullbacks have degree bounds \(d_1,\ldots,d_N\), sampling from a finite field of size at least \((\sum_\ell d_\ell)/\delta\) gives, by \(\mathrm{lem:schwartz\mbox{-}zippel\mbox{-}bound}\) and the union bound,
\[
\Pr(\text{any wrong identity answer})\le\sum_{\ell=1}^N d_\ell/|\mathbb F|\le\delta.
\]
The result follows on the complementary event.
%%% FIELD proof-regular-cell
The supported real source space contains the nonempty Euclidean-open set \(\mathcal S_+=\{(\Lambda,\Omega):\Omega\succ0\}\). A complex polynomial vanishing on a real open box has all coefficients zero by repeated univariate specialization; hence \(\mathcal S_+\) is Zariski dense. On the certificate event, \(\mathrm{thm:exceptional\mbox{-}set\mbox{-}validity}\) makes the source pullback of \(h_{G,e}\) nonzero. Thus some \((\Lambda_0,\Omega_0)\in\mathcal S_+\) has covariance \(\Sigma_0\notin\mathcal E_{G,e}\).

The model image, exceptional complement, discriminant complement, and finitely many lifting sign conditions are semialgebraic. By \(\mathrm{lem:semialgebraic\mbox{-}stratification}\), refine them into smooth strata and take the connected sign stratum containing \(\Sigma_0\). Root number and extension labels are constant there. A sufficiently small closed ball in a smooth chart is a nonempty compact \(\mathcal K_G\) with nonempty relative interior.

For \(\Sigma\in\mathcal K_G\subset\mathcal M_G\), the model definition supplies real supported \((\Lambda,\Omega)\) with \(\Omega\succ0\). Since \(\Omega=(I_n-\Lambda)^{\mathsf T}\Sigma(I_n-\Lambda)\), its missing entries are exactly the residuals \(F_{ij}\), hence vanish. Therefore \(\lambda_e\in\mathcal R_e(G,\Sigma)\), proving nonemptiness.
%%% FIELD proof-root-separation
Every denominator, leading-coefficient factor, and squarefreeness factor is continuous and nonzero on compact \(\mathcal K_G\), so its absolute value has a positive minimum. The coefficients \(c_j(\Sigma)\) of
\[
p_e(\Sigma,t)=t^d+\sum_{j<d}c_j(\Sigma)t^j,\qquad d=d_e(G),
\]
are therefore uniformly bounded, say by \(M\). The monic Cauchy argument confines every root to \(|z|\le1+M\).

The root-incidence set over \(\mathcal K_G\) in this disk is compact. The derivative \(\partial_t p_e\) is nonzero there by squarefreeness and consequently has a positive absolute minimum. If distinct roots had distances tending to zero, compactness would give a limiting double root, contradicting that derivative bound. Thus some \(\gamma>0\) uniformly separates all complex roots and hence all retained real roots.
%%% FIELD proof-regular-inference
Let \(W_r=\operatorname{vech}(X^{(r)}X^{(r)\mathsf T}-\Sigma)\). Differentiating the centered Gaussian moment-generating function four times at zero gives
\[
\operatorname{Cov}(W_{r,ij},W_{r,kl})
=\sigma_{ik}\sigma_{jl}+\sigma_{il}\sigma_{jk}=(V_\Sigma)_{ij,kl}.
\]
For fixed \(t\), Taylor expansion of the characteristic function yields
\[
\mathbb E_\Sigma e^{it^{\mathsf T}W_r/\sqrt m}
=1-\frac{t^{\mathsf T}V_\Sigma t}{2m}+O(m^{-3/2})
\]
uniformly on compact \(\mathcal K_G\), because Gaussian sixth moments are uniformly bounded. Raising to the \(m\)-th power proves the uniform covariance central limit theorem.

Smoothness gives
\[
\pi_G(S_m)-\Sigma=D\pi_G(\Sigma)[S_m-\Sigma]+O(\|S_m-\Sigma\|^2).
\]
Implicit differentiation of \(p_e(\Sigma,r_j(\Sigma))=0\), using the derivative margin, gives
\[
Dr_j(\Sigma)[H]=-\frac{D_\Sigma p_e(\Sigma,r_j)[H]}{\partial_t p_e(\Sigma,r_j)}.
\]
Uniformly bounded derivatives on the compact neighborhood and another Taylor expansion give the stated joint normal limit with covariance
\[
Dr(\Sigma)D\pi_G(\Sigma)V_\Sigma D\pi_G(\Sigma)^{\mathsf T}Dr(\Sigma)^{\mathsf T}.
\]

All certificate factors and root gaps have positive margins. Since \(\|S_m-\Sigma\|=O_{\mathbb P}(m^{-1/2})\) uniformly, with uniform probability tending to one the sample is in the projection tube, every circuit remains valid, and each estimated root lies within \(\gamma/3\) of exactly one population root. Call this event \(\mathcal H_m\). Constant extension labels and \(\mathrm{thm:real\mbox{-}root\mbox{-}lifting}\) then retain exactly the same branches. On \(\mathcal H_m\),
\[
d_H(\widehat{\mathcal R}_{e,m},\mathcal R_e(G,\Sigma))
=\max_j|r_j(\pi_G(S_m))-r_j(\Sigma)|=O_{\mathbb P}(m^{-1/2}).
\]
Also \(\Pr(\widehat{\mathcal R}_{e,m}=\varnothing)\le\Pr(\mathcal H_m^c)\) uniformly.

In smooth local coordinates on the constant-dimensional stratum, the projected covariance is positive definite because \(V_\Sigma\) is positive definite and \(D\pi_G\) is the identity on the tangent space. Its plug-in estimate is uniformly consistent. The Wald ellipsoid at the chi-square \(1-\alpha\) quantile therefore covers the projected population covariance with probability \(1-\alpha+o(1)\), uniformly. Mapping this single ellipsoid through every \(r_j\) simultaneously covers all branches whenever it covers the covariance. This proves the confidence-neighborhood assertion without branch selection.
%%% FIELD proof-singleton-inference
If \(\mathcal R_e(G,\Sigma)=\{r_1(\Sigma)\}\), the bijection on \(\mathcal H_m\) makes \(\widehat{\mathcal R}_{e,m}\) a singleton with uniform probability tending to one. Calling its element \(\widehat r_m\), the Hausdorff bound gives uniform consistency. Taking the first coordinate of the joint limit in \(\mathrm{thm:regular\mbox{-}finite\mbox{-}set\mbox{-}inference}\) gives
\[
\sqrt m\{\widehat r_m-r_1(\Sigma)\}\rightsquigarrow N(0,W_{11}(\Sigma)),
\]
where \(W_{11}\) is the corresponding diagonal entry of the displayed joint covariance; this includes a degenerate normal law when \(W_{11}=0\).
%%% FIELD partial-sharp-frontier
The proved results imply
\[
D_k(n)\le2^{k+1},\qquad
D_0(2)=D_0(3)=1,\qquad D_0(n)=2\ (n\ge4).
\]
Choosing the one-sink family with \(j=\min\{k,\lfloor(n-4)/2\rfloor\}\) excess arrows and appending fully bidirected root-children gives, for \(n\ge4\),
\[
D_k(n)\ge2^{\min\{k+1,\lfloor(n-2)/2\rfloor\}}.
\]
Thus \(D_k(n)=2^{k+1}\) when \(n\ge2k+4\). The exact five-vertex computation below also gives \(D_k(5)\ge4\) for \(k\ge2\). Equality remains unproved for positive \(k\) in the strip \(n<2k+4\).
%%% FIELD partial-five-node
For
\[
D=\{0\to1,0\to2,1\to3,2\to3,1\to4,2\to4\},\qquad
B=\{0\leftrightarrow i:1\le i\le4\},
\]
query \(f=\lambda_{24}\). At the rational source point with directed coefficients \((a,b,c,d,u,f)=(2,3,4,5,6,7)\), disturbance diagonal \((30,31,32,33,34)\), and \(\omega_{0i}=i\), strict diagonal dominance gives \(\Omega\succ0\), and
\[
\Sigma=\begin{pmatrix}
30&61&92&707&1014\\
61&155&187&1561&2247\\
92&187&314&2327&3332\\
707&1561&2327&17981&25747\\
1014&2247&3332&25747&36972
\end{pmatrix}.
\]
The six residuals are
\[
\begin{aligned}
&30ab-92a-61b+187,\\
&61ac+92ad-707a-155c-187d+1561,\\
&61au+92af-1014a-155u-187f+2247,\\
&61bc+92bd-707b-187c-314d+2327,\\
&61bu+92bf-1014b-187u-314f+3332,\\
&155cu+187cf-2247c+187du+314df-3332d-1561u-2327f+25747.
\end{aligned}
\]
Exact rational Buchberger reduction gives the lexicographic basis
\[
\begin{aligned}
0={}&384915840a+425450175590061f^3-8965411803407607f^2\\
&+62974399432228195f-147445028655847225,\\
0={}&741827520b-593375419233f^3+12321625507347f^2\\
&-85271405126447f+196665729339485,\\
0={}&615040c+25798931271f^3-543666188277f^2\\
&+3818873142057f-8941504654939,\\
0={}&4d-3f+1,\\
0={}&153760u+8599643757f^3-181222062759f^2\\
&+1272957714019f-2980501654153,
\end{aligned}
\]
together with
\[
q(f)=(f-7)(8599643757f^3-181150021425f^2
+1271943861499f-2976933819575).
\]
The exact Euclidean algorithm gives \(\gcd(q,q')=1\). The five linear equations recover one value of every other coefficient from each root, so the specialized ideal consists of four reduced points with distinct \(f\)-coordinates. The block graph is \(K_4\) with exponents \((1,1,2,2)\), and
\[
[X_1X_2X_3^2X_4^2]\prod_{i<j}(X_i+X_j)=4.
\]
The multihomogeneous upper bound and the four simple formal branches imply the characteristic-zero identity \(d_{\,2\to4}(G)=4\). This lifts the modular quartic, but the other rooted orbits with \(n\le5\), \(k\le2\) remain uncensused.
%%% FIELD partial-fpt
The quadratic core has at most \(2k\) variables, polynomially many quadratic equations, and polynomially many recorded divisors. Enumerating independence candidates and divisor strata costs \(n^{O(k+1)}\) identity tests, while every zero-dimensional stratum has degree at most \(2^{2k}\). Over \(K_G(U)\), multiplication by \(\lambda_e\) therefore has size at most \(2^{2k}\); its first matrix linear dependence gives the coordinate minimal polynomial. Pure transcendence shows that polynomial is \(p_e\in K_G[t]\), and a denominator-avoiding evaluation of \(U\) removes directed coordinates from its circuits. What remains unproved is one explicit prime-component algorithm that simultaneously handles all divisors, proves covariance-only descent and slice coverage, emits a checkable transcendence certificate, and bounds every intermediate circuit and PIT degree by \(n^{C(k+1)}\).
%%% FIELD partial-real-lifting
For any complete real supported lift, acyclicity makes \(R=I_n-\Lambda\) invertible and
\[
x^{\mathsf T}\Omega x=x^{\mathsf T}R^{\mathsf T}\Sigma Rx=(Rx)^{\mathsf T}\Sigma(Rx)>0
\]
for \(x\ne0\). Hence \(\Omega\succ0\), and the missing residual equations are exactly its required zero entries. Thus a real root is retained exactly when its semialgebraic incidence fiber has a real point. The unresolved part is the finite \(\mathsf{LiftCert}\): arbitrary real covariances have no finite encoding in the frozen arithmetic model, and positive-dimensional pole divisors are not covered by the cited bivariate rational-univariate machinery.
